\section{Non-Chronological Backjumping}

An issue I mentioned earlier with our initial clause learning algorithm is that we backtrack 
only one level per contradiction. Not one level at a time, just one level. The reason why this 
is an issue is that we generally never get to use what we learned from conflict resolution to 
flip certain assignments, which is why we still relied on the DPLL notion of ``trying the other 
branch.'' Non-chronological backtracking, or ``backjumping'' is the new technique we'll use 
to fully exploit the learned clause. 

We only need to make a very minimal interface change. Now that we won't use a ``triedFalse''
variable for our trail starts, our trail starts array can now just be an integer array 
where each element is just the index in the trail where a new decision level starts.
\begin{lstlisting}
    public interface:
        SATSolver(string input)
        bool solveCNF()

    private interface:
        ...
        int[] trailStarts
        ...
\end{lstlisting}

Next, since we're only changing the backtracking logic, we can keep 
everything else the same in our algorithm. Before we move onto code, answer this:
what decision level should be backjump to after we learn a clause? Suppose we learn 
this clause: 

$$ (1 \vee 2 \vee \neg 4) $$

And each literal's decision level is its index (1 was assigned on decision level 1, 
2 was assigned on decision level 2, 4 was assigned on level 4). What decision level 
should be backjump to? Should we backjump to 1? 2? Stay on 4? We want to backjump to 
decision level 2 because that way, every literal in the clause other than $4$ will 
be False, and $4$ will be Unassigned. That way, we can propagate immediately on the 
learned clause! So the rule for backjumping after learning a new clause is to backjump 
to the second largest decision level present in the clause. Here's how we'll do it:

\begin{lstlisting}
bool solve() {
    var[1..numVars] assignment

    for each (var in assignment) {
        var = Unassigned
    }

    if (createWatched(assignment) == False) {
        return False
    }

    int decisionLevel = 0

    if (trail.size > 0) {
        append(trailStarts, {0, true})

        if (propagate(assignment, decisionLevel) == False) {
            return False
        }
    }

    while (True) {
        if (trailHead == trail.size) {
            //Perform an assignment
            ...
        }

        optional int[] contradiction = propagate(assignment)

        if (contradiction has a value) {
            //resolve conflict
            ...

            if (conflict.size == 0) {
                return False
            }

            int maxLevel = decisionLevel
            int secondLargest = -1
            for (i = 0 up to conflict.size) {
                int idx = abs(conflict[i])
                int decisionLvl = trail[var_to_trail[idx]].decisionLevel;
                if (decisionLvl != maxLevel) {
                    secondLargest = max(secondLargest, decisionLvl)
                }
            }

            int backjumpTo = secondLargest

            while (trail.size > 0 && trail.last.decisionLevel > backjumpTo) {
                int decision_start = trailStarts.last
                while (trail.size > decision_start) {
                    int lit = trail.last.lit
                    int idx = abs(lit)
                    assignment[idx] = Unassigned
                    var_to_trail[idx] = -1
                    erase(trail, trail.last)
                }
                erase(trailStarts, trailStarts.last)
            }

            if (trailStarts.size == 0) {
                return False
            }

            decisionLevel = backjumpTo

            trailHead = trailStarts.last

            append(clauses, conflict)
            createWatched(clauses.size)

            continue
        }
    }
}
\end{lstlisting}

The code remains largely the same. However, we now want to check that the learned clause isn't empty
so we can avoid doing large amounts of work. When can the learned clause become empty? To answer 
that question, let's answer this other question: ``what does an empty clause mean?'' From our 
section on Watched Literals, an empty clause means that the equation is unsatisfiable because 
no variable assignment can satisfy that clause. So if we learn an empty clause, that means that 
we're now learning that there is no satisfactory assignment of variables, thus our equation is 
unsatisfiable.

\subsection{Non-Chronological Backjumping - Results}
Okay so I know I said that we would make leaps forward in the next sections.
I may have lied. Here are the results for Non-Chronological Backjumping:

\noindent
uf20.91: \textbf{380ms}, down 45ms from Clause Learning.

\noindent
UUF50.218.1000: \textbf{11291ms}, up 1639ms from Clause Learning.

Why did our time to solve unsatisfiable equations go up so significantly? It has to do 
with the fact of how we're learning. We're not trying every assignment like brute force, but 
we are slowly learning that no matter what the assignment is, our equation is unsatisfiable. 
This takes a lot of backtracking and memory to learn, possibly attributing to our time increase.
It's hard to say for certain.

\begin{FVerbatim}
On all 1000 equations in uf20-91:

Backjumping:               380ms (0.38 seconds).
Clause Learning:           425ms (~0.43 seconds).
Iterative DPLL (AbP):      328ms (~0.33 seconds).
Iterative DPLL (PbA):      324ms (~0.32 seconds).
Watched Literals:          349ms (~0.35 seconds).
Pure Literal Elimination:  1548ms (~1.55 seconds).
Unit Propagation:          868ms (~0.87 seconds).
Recursive Backtracking:    1,540,613ms (~25.68 minutes).
Brute Force:               1,770,851ms (~29.51 minutes).
Short-Circuited (Basic):   49,874ms (~50 seconds).
Parallel:                  10,589ms (~10.5 seconds).

On all 1000 equations in UUF50.218.1000:

Backjumping:              11,291ms (~11.29 seconds).
Clause Learning:          9,652ms (~9.65 seconds).
Iterative DPLL (AbP):     6,369ms (~6.37 seconds).
Iterative DPLL (PbA):     6,365ms (~6.37 seconds).
Watched Literals:         7,130ms (~7.13 seconds).
Pure Literal Elimination: 231,107ms (~3.85 minutes).
Unit Propagation:         140,041ms (~2.34 minutes).
\end{FVerbatim}